\documentclass{article}
\usepackage[T1]{fontenc}
\usepackage{lmodern}
\usepackage{graphicx}
\usepackage{amsmath,amssymb}
\usepackage[a4paper,margin=2cm]{geometry}
\usepackage[absolute,overlay]{textpos}
\setlength{\TPHorizModule}{1mm}
\setlength{\TPVertModule}{1mm}
\usepackage[indonesian]{babel}
\usepackage{hyperref}
\setlength{\emergencystretch}{2em}
% unit: Halaman 20

\begin{document}

% === Halaman 20 (python/native) ===

20

Kelemahan OTP

• Meskipun OTP menawarkan keamanan yang sempurna, tetapi ia tidak umum 

digunakan dalam aplikasi praktis (aplikasi komersil maupun aplikasi lainnya).

• Alasan:

 1. Tidak sangkil, karena panjang kunci =  panjang pesan. 

     Makin panjang pesan, makin besar ukuran kuncinya. Butuh komputasi yang

       berat untuk membangkitkan milyaran  karakter-karakater yang benar-benar

       acak. Untuk pesan yang sangat panjang atau komunikasi yang berlangsung terus-

menerus, kebutuhan kuncinya menjadi sangat besar.

   2. Karena kunci dibangkitkan harus benar-benar acak (truly random), maka ‘tidak 

mungkin’  pengirim dan  penerima membangkitkan kunci acak yang sama secara 

simultan (bersamaan).

\end{document}
